\documentclass[12pt]{article}
\usepackage[utf8]{inputenc}
\usepackage[T1]{fontenc}
\usepackage{hyperref}
\title{Baka Language Resources for K-12 Curriculum Development}
\author{Compiled by Claude Sonnet 5.5 (claude-sonnet-5-5) for Andr\'as Kornai}
\date{}
\begin{document}
\maketitle

\section{Introduction}
Baka (ISO 639-3 \texttt{bkc}) is an Ubangian language of the Baka people of Cameroon and Gabon, written in Latin script. Scholarly literature exists (grammar, nominal morphology, narrative discourse, orthography guide, texts), but no school materials were found.

\section{Primary Educational Resources}
No primers or textbooks were found in the search results.

\section{Linguistic and Reference Materials}
\begin{itemize}
\item Wikipedia overview: \url{https://en.wikipedia.org/wiki/Baka_language}
\item Glottolog record for Kilian-Hatz (1995), \emph{Das Baka}: \url{https://glottolog.org/resource/reference/id/63144}
\item Endangered Languages Project bibliography for Baka (East Region, Cameroon), listing among others Brisson's French--Baka lexicon (1984) and Leonard's work: \url{https://www.endangeredlanguages.com/lang/5989/bibliography}
\item Grambank entry for Baka: \url{https://grambank.clld.org/languages/baka1272}
\end{itemize}

\section{Religious and Community Resources}
No Bible translation or community organisation was identified in the search results.

\section{Digital and Audio Resources}
\begin{itemize}
\item Glosbe online dictionary page (user-contributed, few details): \url{https://glosbe.com/sdb/bkc}
\end{itemize}

\section{Curriculum Development Context}
The existing orthography guide and oral narratives (SIL Cameroon, Leonard 2009, 2011) and the Baka--French dictionary (Brisson and Boursier 1979) are likely the strongest foundations, but they were seen only as bibliographic entries, not as accessible files. Their availability should be checked directly.
\end{document}